% GENERATED FILE. Edit canonical lessons, then run npm run generate:books.
% canonical-source-hash: fnv1a64:dc0b0a891fc0a78d
% canonical-lessons: RU-C440-perekhod, RU-C440-perekhodit, RU-C440-naberezhnaya, RU-C440-bulvar, RU-C440-alleya

\chapter{Crossing the street}
\label{ch:ru-crossing-the-street}
\hlchaptermodality{440}

\begin{chapteropening}
\textbf{By the end of this chapter:} \emph{I can talk about crossings and walks: a crossing, crossing over, the embankment, the boulevard and an avenue of trees.}

Complete the last lesson of chapter 440: I can talk about crossings and walks: a crossing, crossing over, the embankment, the boulevard and an avenue of trees.
\end{chapteropening}

\section[\texorpdfstring{\ru{переход} (perekhód)}{переход (perekhód)}]{\ru{переход} (perekhód) — a crossing; an underpass}

\label{lesson:RU-C440-perekhod}



\begin{quote}
Before the new one: say the Russian for to take a short cut, then the Russian for to go round.
\end{quote}

\subsection*{You'll want to know: \ru{переход}}
\textbf{\ru{переход}} (\emph{perekhód}) — \textquotedblleft{}a crossing; an underpass\textquotedblright{}.

\begin{quote}
\ru{Подземный} \ru{переход} \ru{рядом} \ru{с} \ru{остановкой}.
\end{quote}

\begin{quote}
\emph{Podzémnyy perekhód ryádom s ostanóvkoy.}
\end{quote}

\begin{quote}
The underpass is next to the bus stop.
\end{quote}

A masculine noun. \emph{\ru{Пешеходный} \ru{переход}} — a pedestrian crossing; \emph{\ru{подземный} \ru{переход}} — an underpass.

\begin{cousinweb}
\emph{\ru{переход}} is built from \emph{\ru{пере}-} 'across' and \emph{\ru{ход}} 'a going', which is inherited from Proto-Slavic \emph{xodŭ} 'motion' (EtymDB).
\end{cousinweb}

\begin{grammarlens}[title={What you've built}]
One more word for telling, explaining and describing at B1.
\end{grammarlens}

\subsection*{Your turn}
\begin{itemize}
\raggedright
  \item \emph{Say it:} \emph{perekhód}
  \item \emph{Say it:} \emph{perekhód}, once more
  \item \emph{Say it:} a sentence of your own with \emph{perekhód}
  \item \emph{Recall:} say the Russian for to take a short cut, then the Russian for to go round, then say \emph{perekhód} again
\end{itemize}

\begin{tcolorbox}[breakable,colback=teal!4,colframe=teal!35!black,title={Before you move on}]
What does \ru{переход} mean? (\textquotedblleft{}A crossing; an underpass\textquotedblright{}.) Say it once more.
\end{tcolorbox}
\section[\texorpdfstring{\ru{переходить} (perekhodít)}{переходить (perekhodít)}]{\ru{переходить} (perekhodít) — to cross (a street); to go over}

\label{lesson:RU-C440-perekhodit}



\begin{quote}
Before the new one: say the Russian for to go round, then the Russian for a crossing.
\end{quote}

\subsection*{You'll want to know: \ru{переходить}}
\textbf{\ru{переходить}} (\emph{perekhodít}) — \textquotedblleft{}to cross (a street); to go over\textquotedblright{}.

\begin{quote}
\ru{Дети} \ru{переходят} \ru{улицу} \ru{на} \ru{зелёный} \ru{свет}.
\end{quote}

\begin{quote}
\emph{Déti perekhódyat úlitsu na zelyónyy svet.}
\end{quote}

\begin{quote}
The children cross the street on the green light.
\end{quote}

Imperfective; perfective \emph{\ru{перейти}} (\emph{\ru{перейду}, \ru{перейдёшь}}). Takes the accusative or \emph{\ru{через}}.

\begin{cousinweb}
\emph{\ru{переходить}} is built from \emph{\ru{пере}-} 'across' and \emph{\ru{ходить}} 'to walk'. EtymDB takes the root from Proto-Slavic \emph{xodŭ} 'motion'.
\end{cousinweb}

\begin{grammarlens}[title={What you've built}]
One more word for telling, explaining and describing at B1.
\end{grammarlens}

\subsection*{Your turn}
\begin{itemize}
\raggedright
  \item \emph{Say it:} \emph{perekhodít}
  \item \emph{Say it:} \emph{perekhodít}, once more
  \item \emph{Say it:} a sentence of your own with \emph{perekhodít}
  \item \emph{Recall:} say the Russian for to go round, then the Russian for a crossing, then say \emph{perekhodít} again
\end{itemize}

\begin{tcolorbox}[breakable,colback=teal!4,colframe=teal!35!black,title={Before you move on}]
What does \ru{переходить} mean? (\textquotedblleft{}To cross (a street); to go over\textquotedblright{}.) Say it once more.
\end{tcolorbox}
\section[\texorpdfstring{\ru{набережная} (náberezhnaya)}{набережная (náberezhnaya)}]{\ru{набережная} (náberezhnaya) — an embankment, a riverside walk}

\label{lesson:RU-C440-naberezhnaya}



\begin{quote}
Before the new one: say the Russian for a crossing, then the Russian for to cross.
\end{quote}

\subsection*{You'll want to know: \ru{набережная}}
\textbf{\ru{набережная}} (\emph{náberezhnaya}) — \textquotedblleft{}an embankment, a riverside walk\textquotedblright{}.

\begin{quote}
\ru{Вечером} \ru{мы} \ru{гуляли} \ru{по} \ru{набережной}.
\end{quote}

\begin{quote}
\emph{Vécherom my gulyáli po náberezhnoy.}
\end{quote}

\begin{quote}
In the evening we walked along the embankment.
\end{quote}

An adjective used as a noun; the stress is on the first syllable.

\begin{cousinweb}
\emph{\ru{набережная}} is built from \emph{\ru{на}} 'on' and \emph{\ru{берег}} 'a bank, a shore' with an adjective ending: the street along the bank. \emph{\ru{Берег}} is inherited from Proto-Slavic \emph{bergŭ} (EtymDB).
\end{cousinweb}

\begin{grammarlens}[title={What you've built}]
One more word for telling, explaining and describing at B1.
\end{grammarlens}

\subsection*{Your turn}
\begin{itemize}
\raggedright
  \item \emph{Say it:} \emph{náberezhnaya}
  \item \emph{Say it:} \emph{náberezhnaya}, once more
  \item \emph{Say it:} a sentence of your own with \emph{náberezhnaya}
  \item \emph{Recall:} say the Russian for a crossing, then the Russian for to cross, then say \emph{náberezhnaya} again
\end{itemize}

\begin{tcolorbox}[breakable,colback=teal!4,colframe=teal!35!black,title={Before you move on}]
What does \ru{набережная} mean? (\textquotedblleft{}An embankment, a riverside walk\textquotedblright{}.) Say it once more.
\end{tcolorbox}
\section[\texorpdfstring{\ru{бульвар} (bulvár)}{бульвар (bulvár)}]{\ru{бульвар} (bulvár) — a boulevard}

\label{lesson:RU-C440-bulvar}



\begin{quote}
Before the new one: say the Russian for to cross, then the Russian for an embankment.
\end{quote}

\subsection*{You'll want to know: \ru{бульвар}}
\textbf{\ru{бульвар}} (\emph{bulvár}) — \textquotedblleft{}a boulevard\textquotedblright{}.

\begin{quote}
\ru{Вдоль} \ru{бульвара} \ru{растут} \ru{старые} \ru{деревья}.
\end{quote}

\begin{quote}
\emph{Vdol bulvára rastút stáryye derévya.}
\end{quote}

\begin{quote}
Old trees grow along the boulevard.
\end{quote}

Masculine.

\begin{cousinweb}
\emph{\ru{бульвар}} was borrowed from French \emph{boulevard} (EtymDB), which goes back to Middle Dutch \emph{bollewerc} 'a bulwark, a bastion'.
\end{cousinweb}

\begin{grammarlens}[title={What you've built}]
One more word for telling, explaining and describing at B1.
\end{grammarlens}

\subsection*{Your turn}
\begin{itemize}
\raggedright
  \item \emph{Say it:} \emph{bulvár}
  \item \emph{Say it:} \emph{bulvár}, once more
  \item \emph{Say it:} a sentence of your own with \emph{bulvár}
  \item \emph{Recall:} say the Russian for to cross, then the Russian for an embankment, then say \emph{bulvár} again
\end{itemize}

\begin{tcolorbox}[breakable,colback=teal!4,colframe=teal!35!black,title={Before you move on}]
What does \ru{бульвар} mean? (\textquotedblleft{}A boulevard\textquotedblright{}.) Say it once more.
\end{tcolorbox}
\section[\texorpdfstring{\ru{аллея} (alléya)}{аллея (alléya)}]{\ru{аллея} (alléya) — an avenue (of trees), a path in a park}

\label{lesson:RU-C440-alleya}



\begin{quote}
Before the new one: say the Russian for an embankment, then the Russian for a boulevard.
\end{quote}

\subsection*{You'll want to know: \ru{аллея}}
\textbf{\ru{аллея}} (\emph{alléya}) — \textquotedblleft{}an avenue (of trees), a path in a park\textquotedblright{}.

\begin{quote}
\ru{Мы} \ru{шли} \ru{по} \ru{длинной} \ru{аллее}.
\end{quote}

\begin{quote}
\emph{My shli po dlínnoy alléye.}
\end{quote}

\begin{quote}
We walked along a long avenue of trees.
\end{quote}

Feminine.

\begin{cousinweb}
\emph{\ru{аллея}} is usually traced to French \emph{allée} 'a path', which EtymDB derives from Old French \emph{aler} 'to go'. EtymDB does not list the Russian word.
\end{cousinweb}

\begin{grammarlens}[title={What you've built}]
One more word for telling, explaining and describing at B1.
\end{grammarlens}

\subsection*{Your turn}
\begin{itemize}
\raggedright
  \item \emph{Say it:} \emph{alléya}
  \item \emph{Say it:} \emph{alléya}, once more
  \item \emph{Say it:} a sentence of your own with \emph{alléya}
  \item \emph{Recall:} say the Russian for an embankment, then the Russian for a boulevard, then say \emph{alléya} again
\end{itemize}

\begin{tcolorbox}[breakable,colback=teal!4,colframe=teal!35!black,title={Before you move on}]
What does \ru{аллея} mean? (\textquotedblleft{}An avenue (of trees), a path in a park\textquotedblright{}.) Say it once more.
\end{tcolorbox}
